\section{总结与延伸}

\subsection{各方法适用场景指南}

\begin{table}[H]
\centering
\begin{tabular}{p{0.22\textwidth}p{0.35\textwidth}p{0.33\textwidth}}
\toprule
\textbf{场景} & \textbf{推荐方法} & \textbf{理由} \\
\midrule
图像修复 & Score Manipulation (RePaint) 或 DPS & 掩码操作简单，两类方法均可 \\
超分辨率 & 伪逆引导 & 下采样是线性操作，可利用解析解 \\
高斯去模糊 & 伪逆引导 & 卷积是线性操作 \\
CT/MRI 重建 & 伪逆引导 & 投影操作是线性的 \\
非线性失真恢复 & DPS & 需要处理非线性 $\mathcal{A}$ \\
布局控制生成 & Attention Map 操控 & 需要操控空间对应关系 \\
风格迁移 & Score Manipulation & 涉及跨图像的特征交换 \\
\bottomrule
\end{tabular}
\caption{不同应用场景下推理时引导方法的选择指南}
\end{table}

\subsection{本讲核心内容回顾}

本讲系统介绍了\textbf{推理时引导生成}的基本框架和方法，涵盖以下核心内容：

\begin{enumerate}
    \item \textbf{动机}：利用预训练扩散模型的迭代生成特性，在推理时注入引导信号，无需重新训练即可实现条件生成。
    \item \textbf{Score Manipulation}：直接操控去噪分数/注意力图的实用方法（如 RePaint），无理论保证但广泛有效。
    \item \textbf{逆问题框架}：将条件生成建模为逆问题 $y = \mathcal{A}(x) + n$，提供了系统化的分析工具。
    \item \textbf{贝叶斯分解}：$\nabla_{x_t} \log p(x_t|y) = \nabla_{x_t} \log p(x_t) + \nabla_{x_t} \log p(y|x_t)$，将问题归结为近似似然分数。
    \item \textbf{伪逆引导}：基于高斯近似，适用于线性前向映射，近似精度相对较好。
    \item \textbf{DPS}：基于 Delta 近似，适用于任意可微前向映射，实现极简但近似更粗糙。
\end{enumerate}

\subsection{与下一讲的衔接}

本讲覆盖了三大类推理时引导技术中的前两类。下一讲（Lecture 12: Inference-Time Guided Generation 2）将讨论第三类更高级的技术，可能包括：
\begin{itemize}
    \item 更精确的似然分数近似方法
    \item 基于优化的推理时引导（如 DDRM、MCG 等）
    \item 推理时引导在 3D 生成、视频生成等领域的具体应用
\end{itemize}

\subsection{拓展阅读}

\begin{itemize}
    \item \textbf{RePaint}: Lugmayr et al., ``RePaint: Inpainting using Denoising Diffusion Probabilistic Models,'' CVPR 2022.
    \item \textbf{DPS}: Chung et al., ``Diffusion Posterior Sampling for General Noisy Inverse Problems,'' ICLR 2023.
    \item \textbf{DDRM}: Kawar et al., ``Denoising Diffusion Restoration Models,'' NeurIPS 2022.
    \item \textbf{Classifier Guidance}: Dhariwal \& Nichol, ``Diffusion Models Beat GANs on Image Synthesis,'' NeurIPS 2021.
    \item \textbf{关于逆问题与扩散模型的综述}：讲者在课上推荐查阅 survey paper，了解更多分数操控和逆问题求解的技术。
\end{itemize}

%% --- Fin del contenido --- %%

\end{document}
